%-------------------------------------------------------------------------------
% PUBLIC CV: Weizhi Xue (for kjadams2000.github.io, files/Weizhi_Xue_CV.pdf)
% Derived from the private master CV (resume/master_cv.tex). Public items only:
% published papers, the accepted NeurIPS 2026 paper, public presentations,
% public software. No GPA, fellowship amount, phone number or street address.
% Build: latexmk -pdf public_cv.tex
%-------------------------------------------------------------------------------
\documentclass[letterpaper,10.5pt]{article}

\usepackage[empty]{fullpage}
\usepackage{titlesec}
\usepackage[usenames,dvipsnames]{color}
\usepackage{enumitem}
\usepackage[hidelinks]{hyperref}
\usepackage[english]{babel}
\usepackage{tabularx}
\usepackage[T1]{fontenc}
\usepackage{lmodern}
\input{glyphtounicode}
\pdfgentounicode=1

\hypersetup{
  pdftitle={Weizhi Xue - Curriculum Vitae},
  pdfauthor={Weizhi Xue}
}

\pagestyle{empty}
\addtolength{\oddsidemargin}{-0.55in}
\addtolength{\evensidemargin}{-0.55in}
\addtolength{\textwidth}{1.1in}
\addtolength{\topmargin}{-0.65in}
\addtolength{\textheight}{1.3in}

\urlstyle{same}
\raggedbottom
\raggedright
\setlength{\tabcolsep}{0in}

\titleformat{\section}{\vspace{-3pt}\scshape\raggedright\large\bfseries}{}{0em}{}[\color{black}\titlerule\vspace{-4pt}]

\newcommand{\resumeItem}[1]{\item\small{#1\vspace{-1pt}}}
\newcommand{\resumeSubheading}[4]{
  \vspace{-1pt}\item
    \begin{tabular*}{\textwidth}[t]{l@{\extracolsep{\fill}}r}
      \textbf{#1} & \small #2 \\
      \textit{\small #3} & \textit{\small #4} \\
    \end{tabular*}\vspace{-6pt}
}
\newcommand{\resumeProjectHeading}[2]{
    \item
    \begin{tabular*}{\textwidth}{l@{\extracolsep{\fill}}r}
      \small #1 & \small #2 \\
    \end{tabular*}\vspace{-6pt}
}
\newcommand{\resumePub}[1]{\item\small{#1}\vspace{-2pt}}
\renewcommand\labelitemii{$\vcenter{\hbox{\tiny$\bullet$}}$}
\newcommand{\resumeSubHeadingListStart}{\begin{itemize}[leftmargin=0.0in, label={}]}
\newcommand{\resumeSubHeadingListEnd}{\end{itemize}}
\newcommand{\resumeItemListStart}{\begin{itemize}[leftmargin=0.18in, itemsep=0.5pt, parsep=0pt]}
\newcommand{\resumeItemListEnd}{\end{itemize}\vspace{-4pt}}

%-------------------------------------------------------------------------------
\begin{document}

\begin{center}
    {\Huge \scshape Weizhi Xue} \\ \vspace{2pt}
    \small PhD Candidate in Theoretical Chemistry, The University of Chicago $|$ Chicago, IL \\
    \small \href{mailto:weizhixue43@gmail.com}{weizhixue43@gmail.com} $|$
    \href{https://kjadams2000.github.io}{kjadams2000.github.io} $|$
    \href{https://github.com/KJAdams2000}{github.com/KJAdams2000} $|$
    \href{https://scholar.google.com/citations?user=uxH0pmgAAAAJ}{Google Scholar} \\
    \small \href{https://orcid.org/0009-0004-4904-1433}{ORCID 0009-0004-4904-1433} $|$
    \href{https://www.linkedin.com/in/weizhi-xue-29a46a337/}{linkedin.com/in/weizhi-xue-29a46a337}
\end{center}

%-----------EDUCATION-----------
\section{Education}
  \resumeSubHeadingListStart
    \resumeSubheading
      {The University of Chicago}{Chicago, IL}
      {Ph.D.\ Candidate in Theoretical Chemistry; Advisor: Prof.\ Gregory A.\ Voth}{Sep 2023 -- Expected Jun 2028}
      \resumeItemListStart
        \resumeItem{M.S.\ in Chemistry (Dec 2024).}
      \resumeItemListEnd
    \resumeSubheading
      {Peking University}{Beijing, China}
      {B.S.\ in Chemistry, Minor in Physics; College of Chemistry and Molecular Engineering}{Sep 2019 -- Jul 2023}
      \resumeItemListStart
        \resumeItem{``Chemistry Star'' highest departmental honor degree. B.S.\ thesis advisor: Prof.\ Yiqin Gao.}
      \resumeItemListEnd
  \resumeSubHeadingListEnd

%-----------RESEARCH EXPERIENCE-----------
\section{Research Experience}
  \resumeSubHeadingListStart

    \resumeSubheading
      {Graduate Research Assistant, Voth Group}{Sep 2023 -- Present}
      {The University of Chicago, Department of Chemistry}{Chicago, IL}
      \resumeItemListStart
        \resumeItem{\textbf{Unified latent-variable framework for CG energy-based models:} showed that bottom-up coarse-graining (CG) is maximum-likelihood training of an energy-based model, and that Ultra-Coarse-Graining (observable discrete latents) and virtual-particle CG (unobservable particle latents) are two instances of one latent-variable construction.}
        \resumeItem{\textbf{Self-Consistent Ultra-Coarse-Graining (SC-UCG), first author, \textit{J.\ Chem.\ Phys.}\ 2026:} discrete internal states of each CG molecule are inferred self-consistently by \textbf{graph message passing} at simulation time, with a \textbf{Bethe approximation plus AI-based inference} for state--state correlations, removing the hand-designed collective variables of prior UCG.}
        \resumeItem{Developed \textbf{Multilayer Internal State Consistency (MISC)}, an ML training method derived from relative-entropy (KL) minimization that trains the force field \textbf{without iterative sampling} of intermediate models (under 2 min/epoch on one GPU).}
        \resumeItem{\textbf{Temperature transferability:} trained at a single temperature ($T=4.3$), SC-UCG places a 3D-Ising-class chiral symmetry-breaking transition (racemic $\rightarrow$ D-/L-rich) near $T_c\approx4.4$ and recovers the collective state-switching events below $T_c$.}
        \resumeItem{Implemented SC-UCG and lambda-dynamics UCG as an \textbf{MPI-enabled C++ LAMMPS package} (LAMMPS-UCG-dev) with PyTorch / PyTorch Geometric training workflows (TorchREM, OpenMSCG-dev); mentored 2 summer students on testing and applications (2025).}
        \resumeItem{\textbf{Microtubule plus-end tips, first author, \textit{Biophys.\ J.}\ 2026:} analyzed terabyte-scale all-atom trajectories of $\sim$30-million-atom microtubule lattices, built a bottom-up CG tubulin model (20 sites/monomer), and ran tens-of-$\mu$s CG MD of 8- and 40-layer lattices, showing that intrinsic \textbf{lattice instability} drives protofilament clustering and explaining GTP vs.\ GDP differences.}
        \resumeItem{Co-author on data-driven \textbf{equation-free} acceleration of 21--38-million-atom all-atom MD (\textit{Biophys.\ J.}\ 2025; produced the longitudinal vs.\ lateral interaction analysis, Fig.\ 4) and on all-atom/Brownian-dynamics tip-curvature analysis (\textit{Biophys.\ J.}\ 2026).}
        \resumeItem{\textbf{Enhanced-sampling CG (with C.\ Yang):} free-energy matching of protein binding landscapes (ACS Fall 2026, BPS 2026).}
      \resumeItemListEnd

    \resumeSubheading
      {Collaborative Machine-Learning Research: MixUni (NeurIPS 2026, second author)}{2026}
      {With Q.\ Huang, W.\ Wang, M.\ Cao and I.\ Su, The University of Chicago}{Chicago, IL}
      \resumeItemListStart
        \resumeItem{Multi-task model for liquid-mixture properties in \textbf{battery-electrolyte design}: E(3)-equivariant EGNN encoder, transformer mixture pooling, shared \textbf{sparse mixture-of-experts} head, and a physics-informed \textbf{stretched-Arrhenius} conductivity head ($\log\sigma=\log A - B/T^{\gamma}$).}
        \resumeItem{Ran the \textbf{first physics-head probe experiments}, which established the value of physics-informed readouts and set the paper's direction; led the main physics-head comparison (Arrhenius vs.\ VTF vs.\ stretched-Arrhenius vs.\ free MLP) and ran additional chemically out-of-distribution (Butina-cluster split) seed experiments.}
      \resumeItemListEnd

    \resumeSubheading
      {Visiting Undergraduate Research Intern, Voth Group}{Jul 2022 -- Sep 2022}
      {The University of Chicago}{Remote}
      \resumeItemListStart
        \resumeItem{Co-developed \textbf{K-means clustering coarse-graining (KMC-CG)}, clustering on a position- and fluctuation-mode-aware residue embedding; validated on G-actin, HIV-1 CA pentamer and Arp2/3 (\textit{JCTC} 2023, co-first author).}
      \resumeItemListEnd

    \resumeSubheading
      {Undergraduate Research Intern, Gao Group}{Mar 2022 -- May 2023}
      {Peking University, College of Chemistry and Molecular Engineering}{Beijing, China}
      \resumeItemListStart
        \resumeItem{B.S.\ thesis: computational prediction of ligand-induced conformational changes and their free-energy profiles.}
        \resumeItem{Contributed code to \textbf{SPONGE}, a GPU (CUDA C/C++) MD engine; protein--ligand binding simulations, FEP, and induced-fit docking.}
      \resumeItemListEnd

  \resumeSubHeadingListEnd

%-----------PUBLICATIONS-----------
\section{Publications \normalfont\small\upshape($^\dagger$ equal contribution)}
  \begin{enumerate}[leftmargin=0.22in, itemsep=1pt]
    \resumePub{Q.\ Huang, \textbf{W.\ Xue}, W.\ Wang, M.\ Cao, I.\ Su. ``MixUni: Joint multi-property prediction for chemical mixtures with physics-informed heads.'' \textit{Advances in Neural Information Processing Systems (NeurIPS)} (2026), accepted. \href{https://openreview.net/forum?id=e3HTIyxtsJ}{[OpenReview]} \href{https://github.com/qinanh/MixUni}{[Code]}}
    \resumePub{\textbf{W.\ Xue}, X.\ Shuai, G.\ A.\ Voth. ``Next generation of ultra-coarse-graining: Self-consistent inference of critical internal states.'' \textit{J.\ Chem.\ Phys.}\ \textbf{165}, 174111 (2026). \href{https://doi.org/10.1063/5.0352901}{[DOI]} \href{https://arxiv.org/abs/2608.05388}{[arXiv:2608.05388]}}
    \resumePub{\textbf{W.\ Xue}, J.\ Wu, T.\ Bidone, G.\ A.\ Voth. ``Lattice instability drives formation of protofilament clusters at the microtubule plus-end tips.'' \textit{Biophys.\ J.}\ \textbf{125}, 1516--1531 (2026). \href{https://doi.org/10.1016/j.bpj.2026.02.013}{[DOI]}}
    \resumePub{T.\ Sk\'ora, J.\ Wu, D.\ Beckett, \textbf{W.\ Xue}, G.\ A.\ Voth, T.\ C.\ Bidone. ``On the curvature and relaxation of microtubule plus-end tips.'' \textit{Biophys.\ J.}\ \textbf{125}, 406--419 (2026). \href{https://doi.org/10.1016/j.bpj.2025.07.003}{[DOI]}}
    \resumePub{J.\ Wu, S.\ Dasetty, D.\ Beckett, Y.\ Wang, \textbf{W.\ Xue}, T.\ Sk\'ora, T.\ C.\ Bidone, A.\ L.\ Ferguson, G.\ A.\ Voth. ``Data-driven equation-free dynamics applied to many-protein complexes: The microtubule tip relaxation.'' \textit{Biophys.\ J.}\ \textbf{124}, 3853--3866 (2025). \href{https://doi.org/10.1016/j.bpj.2025.01.009}{[DOI]}}
    \resumePub{J.\ Wu$^\dagger$, \textbf{W.\ Xue}$^\dagger$, G.\ A.\ Voth. ``K-means clustering coarse-graining (KMC-CG): A next generation methodology for determining optimal coarse-grained mappings of large biomolecules.'' \textit{J.\ Chem.\ Theory Comput.}\ \textbf{19}, 8987--8997 (2023). \href{https://doi.org/10.1021/acs.jctc.3c01053}{[DOI]}}
  \end{enumerate}

%-----------PRESENTATIONS-----------
\section{Presentations \normalfont\small\upshape(* presenting author)}
  \begin{itemize}[leftmargin=0.18in, itemsep=0pt]
    \resumeItem{\textbf{W.\ Xue}, X.\ Shuai, Z.\ Zhou, C.\ Yang, G.\ A.\ Voth. ``A unified latent-variable framework for coarse-grained energy-based models.'' Poster, UChicago--Caltech AI+Science Conference, Chicago, IL (Oct 2026).}
    \resumeItem{\textbf{W.\ Xue}*, G.\ A.\ Voth. ``Unified ultra-coarse-graining framework for multistate systems and collective transitions.'' Talk, ACS Fall 2026, Chicago, IL (Aug 2026).}
    \resumeItem{\textbf{W.\ Xue}, Z.\ Zhou*, G.\ A.\ Voth. ``CD-REM: A generalizable and scalable bottom-up coarse-graining framework with virtual sites.'' Talk, ACS Fall 2026, Chicago, IL (Aug 2026).}
    \resumeItem{C.\ Yang, C.\ Kuang, \textbf{W.\ Xue}*, \textit{et al.}, G.\ A.\ Voth. ``Bottom-up coarse-graining via free energy matching.'' Poster, ACS Fall 2026, Chicago, IL (Aug 2026).}
    \resumeItem{\textbf{W.\ Xue}*, J.\ Wu, G.\ A.\ Voth. ``Multiscale modeling for microtubule dynamics: From bending relaxation to hydrolysis kinetics.'' Poster, Biophysical Society Annual Meeting, San Francisco, CA (2026).}
    \resumeItem{C.\ Yang, C.\ Kuang, \textbf{W.\ Xue}, D.\ Vasiliauskas, G.\ A.\ Voth. ``Protein binding landscapes with enhanced sampling-driven coarse-graining.'' Biophysical Society Annual Meeting, San Francisco, CA (2026).}
    \resumeItem{\textbf{W.\ Xue}*, J.\ Wu, G.\ A.\ Voth. ``Bottom-up coarse-graining reveals formation of stable protofilament clusters at microtubule plus-end tips.'' Flash talk and poster, Biophysical Society Annual Meeting, Los Angeles, CA (2025).}
    \resumeItem{J.\ Wu, \textbf{W.\ Xue}, G.\ A.\ Voth. ``K-means clustering coarse-graining (KMC-CG).'' Biophysical Society Annual Meeting (2024).}
  \end{itemize}

%-----------SOFTWARE-----------
\section{Software}
  \resumeSubHeadingListStart
    \resumeProjectHeading
      {\textbf{AceCG} $|$ \emph{Core developer; Python, MPI, Numba, LAMMPS} $|$ \href{https://github.com/AceVolca/AceCG}{github.com/AceVolca/AceCG}}{2025 -- Present}
      \resumeItemListStart
        \resumeItem{Led the architecture rewrite of a \textbf{43k-line} MPI-parallel CG force-field training engine into a 7-layer design (potentials $\rightarrow$ topology $\rightarrow$ I/O/config $\rightarrow$ compute runtime $\rightarrow$ trainers/optimizers $\rightarrow$ scheduler $\rightarrow$ workflows); largest contributor by code volume.}
        \resumeItem{Built MPI CPU parallelism, a multitask scheduler/resource pool, unified topology management, parallel compute backends, YAML config parsing and a \textbf{CLI} (\texttt{acg-fm}, \texttt{acg-rem}, \texttt{acg-cdrem}, \texttt{acg-dsm}, \texttt{acg-trajmap}, \dots).}
        \resumeItem{Implemented denoising score matching, noisy/mixed-noise force matching and noisy REM with batched compute backends; added trajectory and linear force mapping, dihedral B-splines, on-the-fly validation simulations and versioned force-field caching; wrote developer docs and prepared the package for PyPI release.}
      \resumeItemListEnd
    \resumeProjectHeading
      {\textbf{LAMMPS-UCG-dev} $|$ \emph{Author; C++, MPI, LAMMPS} $|$ \href{https://github.com/KJAdams2000/LAMMPS-UCG-dev}{github.com/KJAdams2000/LAMMPS-UCG-dev}}{2024 -- 2025}
      \resumeItemListStart
        \resumeItem{LAMMPS package and patcher for lambda-dynamics UCG and SC-UCG, with internal-state inference at simulation time.}
      \resumeItemListEnd
    \resumeProjectHeading
      {\textbf{TorchREM / OpenMSCG-dev} $|$ \emph{Author; PyTorch} $|$ \href{https://github.com/KJAdams2000/TorchREM}{github.com/KJAdams2000/TorchREM}}{2024 -- 2025}
      \resumeItemListStart
        \resumeItem{PyTorch relative-entropy trainer for CG MD; OpenMSCG extensions (UCG-weighted REM, state I/O, Adam).}
      \resumeItemListEnd
  \resumeSubHeadingListEnd

%-----------AWARDS-----------
\section{Honors and Awards}
  \begin{itemize}[leftmargin=0.18in, itemsep=0pt]
    \resumeItem{George Herbert Jones Fellowship, UChicago Chemistry (2025); Olshensky Student Travel Award, UChicago Chemistry (2024, 2025)}
    \resumeItem{``Chemistry Star'' highest honor degree, Peking University (2023); Model Student of Academic Records (2020--2022) with Zheng Yongxi, Peking University and POSCO Scholarships}
    \resumeItem{Gold Medal, Chinese Chemistry Olympiad (2018)}
  \end{itemize}

%-----------TEACHING-----------
\section{Teaching and Mentoring}
  \begin{itemize}[leftmargin=0.18in, itemsep=0pt]
    \resumeItem{Tutor and grader, University of Chicago: CHEM 364 Advanced Statistical Mechanics (Winter 2025, Winter 2026); CHEM 365 Chemical Dynamics (Spring 2025); CHEM 261 Introductory Quantum Mechanics (Fall 2025)}
    \resumeItem{Teaching Assistant, CHEM 111/112/113 General Chemistry sequence, University of Chicago (Fall 2023 -- Spring 2024)}
    \resumeItem{Research mentor, PKU--UChicago Summer Research Program (2024, 2025)}
  \end{itemize}

%-----------SKILLS-----------
\section{Technical Skills}
 \begin{itemize}[leftmargin=0.0in, label={}]
    \small{\item{
     \textbf{Languages \& HPC}{: Python, C/C++, CUDA, MPI (mpi4py), Numba, Bash, SLURM clusters, Git, \LaTeX} \\
     \textbf{ML}{: PyTorch, PyTorch Geometric, NumPy/SciPy; equivariant GNNs, energy-based/latent-variable models, score matching, mixture-of-experts} \\
     \textbf{Simulation}{: LAMMPS, GROMACS, OpenMM, PLUMED, AmberTools, SPONGE, Gaussian, MDAnalysis, OpenMSCG; enhanced sampling, FEP, CG (FM, REM, IBI, UCG)}
    }}
 \end{itemize}

\end{document}
